\section{Constitutive operator and reconstruction of the vacuum channels}
\label{iii:sec:constitutiva}

The constitutive quantities are obtained by an operation on angular
transport. The starting information comprises the channels, their
projectors, and the orientation distinguishing them. This organization makes it possible
to follow separately the construction of the operator, the evaluation of its
spectral functions, and its subsequent dimensional realization.

\subsection{Positive response on two sheets}

In the oriented chamber $x>y>0$ of Section~\ref{iii:sec:hojas}, let
$P_\pm=(I\pm\mathcal R)/2$ and
$T=q_+P_++q_-P_-$, with $0<q_+<q_-<1$. The incidences constructed by the
calendar fix the sectors $90$ and $120$. The constitutive response is
\begin{equation}
 \mathcal V=(I-T^{90})(I-T^{120})^{-1}.
 \label{iii:eq:vacuum-op}
\end{equation}
This definition receives the transport arising from the TPK and preserves the
sheet assignments. The response rule and its domain are an explicit part
of the construction.

The quotient arises from comparing the two calendar incidences on
the same transport, not from fitting four quantities separately. If
$S=T^{30}$ and $\Sigma_j(S)=I+S+\cdots+S^{j-1}$, factoring
$I-S^j$ gives
\begin{equation}
 \mathcal V\,\Sigma_4(S)=\Sigma_3(S),\qquad
 \mathcal V=\Sigma_3(S)\Sigma_4(S)^{-1}.
 \label{iii:eq:completion-response}
\end{equation}
The factors $\Sigma_3$ and $\Sigma_4$ express, respectively, the
sectors $90=3\cdot30$ and $120=4\cdot30$. Since the spectrum of $S$ is
contained in $(0,1)$, $\Sigma_4(S)$ is invertible: the completion
equation determines a unique response for the fixed transport.
This uniqueness concerns the declared constitutive rule and its
incidences. Electromagnetic identification adds the oriented
assignment of sheets and the dimensional lines specified
below.

\begin{proposition}[Decomposition and positivity]
The response is well-defined and
\begin{equation}
 \mathcal V=r_+P_++r_-P_-,\qquad
 r_\pm=\frac{1-q_\pm^{90}}{1-q_\pm^{120}}.
 \label{iii:eq:rchannels}
\end{equation}
Its eigenvalues satisfy $3/4<r_-<r_+<1$ in this chamber.
\end{proposition}
\begin{proof}
Each coefficient $1-q_\pm^{120}$ is positive. The inverse of $I-T^{120}$
is the sum of the projectors divided by those coefficients, proving
\eqref{iii:eq:rchannels}. With $s=q^{30}$, one obtains
\begin{equation}
 f(s)=\frac{1+s+s^2}{1+s+s^2+s^3},\qquad
 f'(s)=-\frac{s^2(3+2s+s^2)}{(1+s+s^2+s^3)^2}<0.
 \label{iii:eq:monotone}
\end{equation}
The limits at $0$ and $1$ are $1$ and $3/4$, respectively. Monotonicity
and the channel ordering complete the proof.
\end{proof}

\subsection{Four coordinates of a common response}

Define the normalized evaluations
\begin{equation}
 \widehat\varepsilon=r_+^2,\qquad
 \widehat\mu=r_-^2,\qquad
 \widehat Z=\frac{r_-}{r_+},\qquad
 \widehat c=\frac1{r_+r_-}.
 \label{iii:eq:four}
\end{equation}
The first pair preserves the quadratic responses assigned to each sheet.
The product captures the common mode; the quotient distinguishes their relative orientation.
The following hold exactly:
\begin{equation}
 \widehat\mu\widehat\varepsilon=\widehat c^{-2},\qquad
 \frac{\widehat\mu}{\widehat\varepsilon}=\widehat Z^2,
 \qquad
 r_+=\frac1{\sqrt{\widehat Z\widehat c}},\quad
 r_-=\sqrt{\frac{\widehat Z}{\widehat c}}.
 \label{iii:eq:four-identities}
\end{equation}
The positive roots fix the inversion. Under sheet conjugation,
read relative to the fixed markings $P_\pm$,
$\widehat\varepsilon$ and $\widehat\mu$ are interchanged,
$\widehat Z$ is inverted, and $\widehat c$ is preserved. The response is then transported
to the chamber $x>-y>0$: the formulas remain valid, with the orderings
$q_-<q_+$ and $r_+<r_-$ reversed. Consequently, common mode
and orientation are data recoverable through complementary observables.

\begin{proposition}[Determination of the quadratic evaluations]
Let $W=r_+P_++r_-P_-$ be a positive response on the two sheets.
With the product and quotient readouts fixed,
\[
 \widehat c^{-1}=\det W=r_+r_-,
 \qquad \widehat Z=r_-/r_+,
\]
there exists a unique positive pair $(\widehat\varepsilon,\widehat\mu)$
satisfying
$\widehat\mu\widehat\varepsilon=\widehat c^{-2}$ and
$\widehat\mu/\widehat\varepsilon=\widehat Z^2$.
It is precisely $(r_+^2,r_-^2)$.
\end{proposition}
\begin{proof}
Positivity permits taking roots unambiguously:
\[
 \widehat\mu=\frac{\widehat Z}{\widehat c}
 =\frac{r_-}{r_+}r_+r_-=r_-^2,\qquad
 \widehat\varepsilon=\frac1{\widehat Z\widehat c}
 =\frac{r_+}{r_-}r_+r_-=r_+^2.
\]
These expressions satisfy both relations and are their only positive
roots.
\end{proof}

Thus the squares are not two independent choices added
to the channels: they are determined by the multiplicative and oriented
readouts once the constitutive relations are fixed. This proposition
establishes uniqueness within that system of rules. Interpreting
the rules as an electromagnetic response retains its additional
physical-representation content.

The determinantal readout of transport and the constitutive response
must also preserve their order. For $T=e^{-xI-y\mathcal R}$, one has
$\det T=e^{-2x}$, whereas
\[
 \det\mathcal V=f(q_+^{30})f(q_-^{30})=\widehat c^{-1}.
\]
Both arise from the same two-sheet operator but evaluate different
functions of its spectrum. In general, applying the determinant first
and then a rational function is not equivalent to taking the determinant of
that function of the operator. Keeping both eigenvalues until the
response is completed preserves orientation information that an isolated
determinant no longer distinguishes.

\subsection{Cubic inversion and angular recovery}

\begin{theorem}[Constitutive inversion in the contractive domain]
For each $r\in(3/4,1)$, there exists a unique $s\in(0,1)$ such that $f(s)=r$.
This $s$ is the unique root in that interval of
\begin{equation}
 r s^3+(r-1)(1+s+s^2)=0.
 \label{iii:eq:cubic}
\end{equation}
The labeled response reconstructs $s_\pm$, $q_\pm=s_\pm^{1/30}$, and
\begin{equation}
 x=-\frac12\log(q_+q_-),\qquad
 y=\frac12\log\frac{q_-}{q_+}.
 \label{iii:eq:recover-angular}
\end{equation}
\end{theorem}
\begin{proof}
The derivative and limits in \eqref{iii:eq:monotone} prove that $f$ is a
decreasing bijection between the stated intervals. Multiplying $f(s)=r$
by its denominator produces \eqref{iii:eq:cubic}. The positive root of order
$30$ recovers the channel. Finally, adding and subtracting
$\log q_\pm=-x\mp y$ gives \eqref{iii:eq:recover-angular}.
\end{proof}

Reconstruction preserves the sheet labels. The complete operator
$\mathcal V$ also preserves the projectors, and the inverse function is applied
by spectral calculus. Recovering two angular coordinates is an
operation defined on this representation; the complete historical
TPK trace preserves additional data.

The interval $r\in(3/4,1)$ and the inverse function of this theorem belong
to the internal normalization of \eqref{iii:eq:four}. If coordinates
expressed in another unit chart are available, first recover
this normalization through the transformation in
Equation~\eqref{iii:eq:dimensions}; then apply cubic inversion.

\subsection{Logarithmic coordinates and refinement}

Let $\mathcal E=\log(r_+r_-)$ and $\mathcal B=\log(r_-/r_+)$. Then
\begin{equation}
 \log\mathcal V=\tfrac12\mathcal E I-\tfrac12\mathcal B\mathcal R.
 \label{iii:eq:logvac}
\end{equation}
This decomposition displays the invariant and oriented components.
In the amplification $\mathcal V_m=\mathcal V\otimes I_{9^m}$, normalized
partial expectations satisfy
$E_m(\mathcal V_{m+1})=\mathcal V_m$. For every polynomial $p$,
$p(\mathcal V_m)=p(\mathcal V)\otimes I_{9^m}$; hence
$E_m(p(\mathcal V_{m+1}))=p(\mathcal V_m)$. Uniform continuity
extends the identity to continuous functions on the spectrum.
Thus projectors, powers, and the logarithm are preserved in the stated tower.

\subsection{Dimensional types of the four evaluations}

With the action, charge, and velocity lines, equipped with sections
$u_S,u_Q,u_C$, the dimensional realization reads
\begin{align}
 Z_{\rm vac}&=\widehat Z u_Su_Q^{-2},&
 c_{\rm vac}&=\widehat c u_C,\\
 \mu_{\rm vac}&=\widehat\mu u_Su_C^{-1}u_Q^{-2},&
 \varepsilon_{\rm vac}&=\widehat\varepsilon u_S^{-1}u_C^{-1}u_Q^2.
 \label{iii:eq:dimensions}
\end{align}
Multiplying and dividing these expressions preserves
$\mu_{\rm vac}\varepsilon_{\rm vac}=c_{\rm vac}^{-2}$ and
$\mu_{\rm vac}/\varepsilon_{\rm vac}=Z_{\rm vac}^2$.
The SI chart subsequently evaluates the sections. Physical comparison requires
presenting their construction and normalization alongside evaluation; the four
normalized coordinates in \eqref{iii:eq:four} retain their dimensionless type.
